% Análise da execução de 18/11/2026 (data/eval/resultados.jsonl). Se a avaliação for refeita,
% revise os números citados aqui com base em resultados.tex.
\subsection{Discussão}

O GraphRAG híbrido teve o melhor acerto nas 27 perguntas respondíveis (83,3\%), contra 74,1\% do RAG vetorial
e 74,1\% do planejador KAG. A vantagem apareceu sobretudo nas perguntas de um salto (90\% contra 75\%) e nas
globais (60\% contra 50\%). Nas multi-hop, o híbrido e o KAG empataram em 87,5\%, pouco acima do vetorial
(83,3\%). A diferença menor nas multi-hop tem uma explicação no próprio corpus: ele foi dividido em um trecho
por artigo, e muitos artigos já citam o dispositivo relacionado. O texto do Tema 987, por exemplo, transcreve o
art.~19 do Marco Civil, e a busca vetorial encontra a resposta sem atravessar o grafo.

O grafo também reduziu as recusas indevidas. O vetorial recusou 11,1\% das perguntas respondíveis, o híbrido
7,4\% e o KAG 3,7\%. Com mais contexto, o gerador encontra base para responder com mais frequência. O custo é
um contexto mais ruidoso: a alucinação bruta, isto é, as afirmações geradas que o verificador rejeitou, ficou
em 20,7\% no híbrido, contra 18,7\% no vetorial e 16,2\% no KAG. O KAG respondeu mais, mas parte dessas respostas
foi incompleta. Em duas perguntas de um salto ele descreveu o que o contexto não trazia em vez de responder, o
que o juiz classificou como incorreto.

O mecanismo de anti-alucinação funcionou nos três modos. Todas as perguntas impossíveis foram recusadas, e
nenhuma afirmação sem sustentação chegou à resposta final, porque o verificador as remove. Os erros que restaram
são recusas ou respostas incompletas, e não respostas inventadas. Uma pergunta multi-hop (a remissão da LGPD à
Lei de Arbitragem) foi recusada pelos três modos, porque a busca não trouxe o inciso que a responde.

O modo híbrido custou cerca de 3,2 vezes mais por pergunta que o vetorial (US\$ 0,0126 contra US\$ 0,0039), por
causa do contexto maior e da consulta Cypher, e foi o mais lento (8,0 s em média, contra 4,5 s). A maior parte da
latência vem das chamadas ao LLM (Cypher, geração e verificação).

\textbf{Limitações.} O conjunto tem 32 perguntas, escritas pela própria equipe a partir do corpus. O juiz é o
mesmo modelo que gera as respostas. Cada configuração rodou uma única vez, sem medir a variância entre execuções,
e respostas a uma mesma pergunta variam entre execuções. A diferença de nove pontos entre o híbrido e o vetorial
corresponde a poucas perguntas e deve ser confirmada com um conjunto maior. Por fim, o modelo de embeddings é
pequeno, e perguntas muito coloquiais às vezes não recuperam o artigo certo.
